\section{Đánh giá và hạn chế}

\subsection{Đối chiếu với mục tiêu}

\begin{rtable}
\begin{tabular}{@{}p{0.26\textwidth} p{0.10\textwidth} p{0.52\textwidth}@{}}
\toprule
\textbf{Yêu cầu} & \textbf{Kết quả} & \textbf{Bằng chứng} \\
\midrule
Chỉ trả lời dựa trên văn bản & Đạt & Prompt ràng buộc + soát trích dẫn tự động; hình~\ref{fig:count} cho thấy đường từ chối hoạt động \\
Truy vết được & Đạt & Trích dẫn bấm được (hình~\ref{fig:chat1}), mở trích đoạn gốc (hình~\ref{fig:source}), mở PDF đúng trang \\
Biết nói “không biết” & Đạt & Ngưỡng liên quan theo từng trích đoạn + câu trả lời cố định \\
Chạy tại chỗ & Đạt & Ollama + Qdrant trong Docker; chế độ cloud là tuỳ chọn, embedding luôn local \\
Chất lượng đo được & Đạt & hit@5 = 1,000 và MRR = 0,967 trên 30 câu hỏi vàng (bảng~\ref{tab:eval}) \\
\bottomrule
\end{tabular}
\captionof{table}{Đối chiếu kết quả với yêu cầu ban đầu.}
\label{tab:doichieu}
\end{rtable}

\subsection{Hạn chế}

\begin{itemize}
  \item \textbf{Bộ câu hỏi vàng do nhóm tự soạn.} 30 câu là nhỏ, và quan trọng hơn, chúng được
  viết bởi người đã đọc văn bản — nên có thiên lệch về từ ngữ so với câu hỏi thật của sinh
  viên. Con số hit@5 = 1,000 vì vậy nên đọc là “không còn lỗ hổng rõ ràng nào trên tập này”,
  chứ không phải “hệ thống không bao giờ trượt”.
  \item \textbf{Chưa đánh giá chất lượng câu trả lời.} Harness đo truy xuất, không đo phần
  sinh. Việc mô hình có diễn giải trung thực trích đoạn hay không hiện chỉ được kiểm bằng mắt
  và bởi cơ chế soát trích dẫn.
  \item \textbf{Phụ thuộc bước rà soát thủ công.} Chất lượng toàn hệ thống bị chặn trên bởi
  chất lượng bản Markdown sau OCR. Với văn bản 57 trang thì chấp nhận được; mở rộng ra hàng
  chục văn bản thì đây là nút thắt.
  \item \textbf{Một văn bản duy nhất.} Pipeline hỗ trợ đa tài liệu nhưng chưa được kiểm chứng ở
  quy mô đó. Khi có nhiều văn bản chồng lấn (quy định cũ – mới), sẽ phát sinh bài toán mới: ưu
  tiên văn bản còn hiệu lực.
  \item \textbf{Endpoint \code{/ingest} chưa xác thực.} Ai truy cập được cổng đều kích hoạt
  được việc nạp lại dữ liệu. Chấp nhận được khi chạy localhost, phải khoá lại nếu triển khai
  thật.
  \item \textbf{\code{/chat} là endpoint đồng bộ}, chạy trong threadpool của FastAPI. Với vài
  chục người dùng đồng thời cần chuyển sang I/O bất đồng bộ.
\end{itemize}

\subsection{Bài học}

Bốn lỗi đáng nhớ nhất trong quá trình làm, đều cùng một kiểu: \emph{trông thì đúng}.

\begin{enumerate}
  \item \textbf{Chuẩn hoá min–max trong bộ rerank.} Nó ánh xạ ứng viên tốt nhất về 1 và kém
  nhất về 0, khiến hai ứng viên chênh nhau 0,01 trông xa nhau như 1,0 và 0,0, và mọi tín hiệu
  từ vựng trở nên vô dụng. Phát hiện nhờ một unit test, không phải nhờ đọc code.
  \item \textbf{Cướp câu hỏi ở bước viết lại} (mục~\ref{sec:rewrite}, hình~\ref{fig:bug}). Hệ
  thống trả lời sai chủ đề nhưng với trích dẫn thật — không một chỉ số kỹ thuật nào bắt được,
  chỉ phát hiện khi có người gõ một câu hỏi vu vơ.
  \item \textbf{Ngưỡng liên quan dùng như một cánh cổng chung} (mục~\ref{sec:nguong}). Chỉ cần
  một trích đoạn đúng là cả năm trích đoạn top-k được giữ theo, nên mọi câu trả lời đều kèm
  đúng từng ấy “nguồn tham khảo” — kể cả câu trả lời kết luận là không tìm thấy gì. Không test
  nào bắt được, vì hệ thống chạy đúng như code mô tả; chỉ người dùng để ý thấy con số không bao
  giờ đổi mới phát hiện ra.
  \item \textbf{Tên file tiếng Việt trong HTTP header.} Header HTTP chỉ mã hoá được latin-1;
  tên file tiếng Việt làm máy chủ trả lỗi 500 khi mở PDF. Kèm theo đó là một lỗi tiềm ẩn: macOS
  lưu tên file ở dạng NFD còn siêu dữ liệu ghi dạng NFC — hai chuỗi như nhau trên macOS nhưng
  khác nhau trong container Linux.
\end{enumerate}

Mẫu số chung: lỗi nguy hiểm nhất của một hệ thống RAG không phải lỗi làm nó hỏng, mà lỗi làm nó
trả lời \textbf{tự tin và sai}, hoặc trình bày thứ trông như bằng chứng mà thực ra không phải.
Đó là lý do đồ án đầu tư vào ba lớp phòng thủ — ngưỡng liên quan theo từng trích đoạn, soát
trích dẫn, và bộ bảo vệ bước viết lại — thay vì chỉ tối ưu điểm số.

\chapterbreak
